\pdfvariable suppressoptionalinfo 611
\providecommand{\CRevisionRound}{2}
\documentclass[10pt]{article}
\usepackage[a4paper,margin=20mm]{geometry}
\usepackage{amsmath,amssymb,amsthm,booktabs,microtype}
\usepackage[T1]{fontenc}
\usepackage{lmodern}
\usepackage[hidelinks]{hyperref}
\newtheorem{theorem}{Theorem}
\newtheorem{proposition}{Proposition}
\newcommand{\diab}{\mathrm{diab}}
\newcommand{\ad}{\mathrm{ad}}
\title{Landau--Zener--Weber Scattering: Exact Connection Data and\newline Finite-Window Unitary Controls}
\author{HCS Research Program}
\date{29 August 2026\quad (revision \CRevisionRound)}
\begin{document}
\maketitle

\begin{abstract}
We analyze the genuinely nonautonomous two-level crossing
$i\dot\psi=[(vt/2)\sigma_z+g\sigma_x]\psi$, $v>0$.  Eliminating one
component gives a pair of parabolic-cylinder (Weber) equations.  Their
connection formula yields, in a fixed diabatic gauge, the exact asymptotic
survival law $P_\diab=\exp(-2\pi g^2/v)$ and the
Stokes/Gamma phase
$\phi_S=\pi/4+\delta(\log\delta-1)+\arg\Gamma(1-i\delta)$.
Wronskian conservation proves an $SU(2)$ scattering matrix, while elementary
calculus gives strict monotonicity and sudden/adiabatic limits.  A separate
80-digit, fixed-step RK4 ledger checks finite windows without calling them an
exact finite-time formula.  The source-local theorem is a unitary/scattering
candidate only: it supplies no arithmetic owner, target divisor, or
Hilbert--Polya bridge.
\end{abstract}

\section{Crossing and conventions}
Set $\hbar=1$, $\sigma_z=\operatorname{diag}(1,-1)$, and use the diabatic
basis at both ends of time:
\begin{equation}
 H(t)=\frac{vt}{2}\sigma_z+g\sigma_x,
 \qquad i\frac{d}{dt}\binom{a}{b}=H(t)\binom{a}{b},
 \qquad v>0,\ g\in\mathbb R .                         \label{eq:model}
\end{equation}
The physical clock is $t$; there is no fitted repetition or target clock.  We
write $\delta=g^2/v$ and fix the asymptotic phase convention by the matrix in
Theorem~\ref{thm:scattering}.  Conjugating by $\sigma_z$ maps $g$ to $-g$.

\begin{theorem}[Weber reduction and scattering law]\label{thm:scattering}
For $g\ne0$, the components of a solution of \eqref{eq:model} satisfy
\begin{align}
 a''+\left(g^2+\frac{v^2t^2}{4}+\frac{iv}{2}\right)a&=0,\\
 b''+\left(g^2+\frac{v^2t^2}{4}-\frac{iv}{2}\right)b&=0.               \label{eq:weber}
\end{align}
After the usual rotation of $t$ and affine rescaling these are Weber
parabolic-cylinder equations.  Unit-flux connection data from $t=-\infty$ to
$t=+\infty$ are
\begin{equation}
 S(\delta)=\begin{pmatrix}
 \sqrt P&-\sqrt{1-P}\,e^{i\phi_S}\\
 \sqrt{1-P}\,e^{-i\phi_S}&\sqrt P
 \end{pmatrix},\quad
 P=e^{-2\pi\delta},\quad
 \phi_S=\frac\pi4+\delta(\log\delta-1)+\arg\Gamma(1-i\delta),       \label{eq:S}
\end{equation}
with $\phi_S(0)=\pi/4$ by continuity.  A different diagonal asymptotic
phase convention conjugates $S$ by diagonal unitaries and leaves $P$ intact.
\end{theorem}

\begin{proof}
From the first component equation,
$b=(ia'-(vt/2)a)/g$.  Substitution into the second gives the first equation
in \eqref{eq:weber}; exchanging components gives the second.  The rotated
equation is the standard parabolic-cylinder equation.  Its connection identity
and constant Wronskian produce the two entries in \eqref{eq:S}; the modulus
identity $|\Gamma(1+i\delta)|^2=\pi\delta/\sinh(\pi\delta)$ reduces the
diagonal coefficient to $e^{-\pi\delta}$.  This is a source-local special
function calculation, with no external spectral data.
\end{proof}

\begin{proposition}[Unitary and monotone invariants]\label{prop:inv}
For real parameters, $H(t)=H(t)^*$ and every finite-time propagator preserves
the Hermitian norm.  The Weber Wronskian gives the same flux identity at the
two ends, so $S^*S=I$ and $\det S=1$.  Moreover
\begin{equation}
 \frac{dP}{d\delta}=-2\pi e^{-2\pi\delta}<0\quad(\delta>0),\qquad
 \left.\frac{\partial P}{\partial g}\right|_v=-\frac{4\pi g}{v}e^{-2\pi g^2/v},
 \label{eq:deriv}
\end{equation}
and $\phi_S'(\delta)=\log\delta-\Re\psi(1-i\delta)$.
\end{proposition}

\begin{proof}
Hermiticity gives $(\psi^*\psi)'=0$; the Wronskian is the equivalent
constant flux in the Weber basis.  Direct multiplication of \eqref{eq:S}
proves the $SU(2)$ identities.  Differentiating $e^{-2\pi\delta}$ gives
\eqref{eq:deriv}; differentiating $\log\Gamma(1-i\delta)$ gives the phase
derivative.
\end{proof}

\clearpage
\section{Limits, boundaries, and a finite-window control}

The crossing parameter organizes the singular faces.  As $\delta\downarrow0$,
\begin{equation}
 P=1-2\pi\delta+O(\delta^2),\qquad 1-P=2\pi\delta+O(\delta^2),
\end{equation}
so the sudden limit $v\to\infty$ is diabatic.  At fixed $g\ne0$, $v\downarrow0$
gives $P\to0$ (adiabatic following).  At $g=0$ the two channels decouple
exactly, and the continuous phase convention is used.  The sign change
$g\mapsto-g$ is a constant $\sigma_z$ gauge, not a new probability law.  The
point $t=0$ is the sole diabatic gap closing; it is the turning point of the
Weber connection problem.

For validation, we integrated both basis vectors of \eqref{eq:model} on
$[-T,T]$ with a fixed 2048-step classical RK4 scheme at 80 decimal digits.
The ledger records all four complex entries, $P_{\rm window}=|U_{11}|^2$,
$|P_{\rm window}-P|$, and $\|U^*U-I\|_{\max}$.  It is a controlled numerical
window, not an exact finite-time propagator.  Representative exact rows are:

\begin{table}[h]
\centering
\small
\begin{tabular}{@{}lrrrcc@{}}
\toprule
case & $v$ & $g$ & $\delta$ & $P_\diab$ & $\phi_S$\\
\midrule
reference & 1 & $1/2$ & $1/4$ & 0.2078795764 & 0.3270619133\\
fast/weak & 4 & $1/3$ & $1/36$ & 0.8398492016 & 0.6741033767\\
slow/strong & $1/4$ & $3/4$ & $9/4$ & 0.0000007249 & 0.0372977224\\
negative $g$ & 2 & $-1/2$ & $1/8$ & 0.4559381278 & 0.4718436007\\
uncoupled & $3/2$ & 0 & 0 & 1 & 0.7853981634\\
\bottomrule
\end{tabular}
\caption{Exact source-local scattering rows in the fixed gauge.}
\label{tab:scattering}
\end{table}

The finite ledger has 15 rows ($T=2,4,8$ for each case).  The largest Gram
residual is $1.56\times10^{-5}$ with the 2048-step control; discrepancies are
reported as finite-window effects.  A producer-independent checker recomputes
the ODE, while a separate SymPy program checks the signs in \eqref{eq:weber},
the Pauli square, $SU(2)$ determinant, monotonicity, and the coupling gauge.
The machine-readable field is written literally as \texttt{P\_diabatic}; the
scope lock is \texttt{NO\_BAD\_EULER\_OR\_ROOT\_NUMBER}.

\section{Relation to the Route-A decision}

The theorem is a complete physical scattering result, but a single swept
crossing has no primitive periodic-orbit repetition law.  The parameter
$\delta$ is a source coupling ratio, not an arithmetic label; $P$ and
$\phi_S$ are not target zeros or divisor data.  Thus the strict tuple is
\begin{center}
\ttfamily (A0\_FAIL, A1\_FAIL, A2\_FAIL, A3\_FAIL,
A4\_UNITARY\_OR\_SCATTERING\_CANDIDATE),
\end{center}
with \texttt{overall=ROUTE\_A\_REJECTED} and
\texttt{route\_b\_invocation\_allowed=false}.  This nonautonomous model is
intentionally distinct from the autonomous Jaynes--Cummings excitation blocks
of C223.  No target operator, Euler factor, root number, automorphy claim, or
Hilbert--Polya bridge is made.

\clearpage
\section*{Reproducibility and source note}

The receipt is generated by
\texttt{code/c224\_landau\_zener\_producer.py} using exact rational sentinels,
80-digit Gamma evaluation, and deterministic RK4.  The independent checker,
symbolic cross-check, clean replay, and hostile mutation suite are run in
separate processes.  The fixed-epoch LuaLaTeX build is repeated twice; the
three revision PDFs are content-distinct and the final \texttt{main.pdf} is
byte-identical to round 2.  The release manifest closes exactly 27 payload
files and excludes itself and build sidecars.

The source audit records the following DOI-verified references:

\begin{thebibliography}{9}
\bibitem{zener1932} C. Zener, ``Non-adiabatic crossing of energy levels,''
\emph{Proceedings of the Royal Society A} 137 (1932), 696--702.
DOI: \href{https://doi.org/10.1098/rspa.1932.0165}{10.1098/rspa.1932.0165}.
\bibitem{vitanov1996} N. V. Vitanov and B. M. Garraway, ``Landau-Zener model:
Effects of finite coupling duration,'' \emph{Physical Review A} 53 (1996),
4288--4304. DOI: \href{https://doi.org/10.1103/PhysRevA.53.4288}{10.1103/PhysRevA.53.4288}.
\bibitem{shevchenko2010} S. N. Shevchenko, S. Ashhab, and F. Nori,
``Landau-Zener-Stuckelberg interferometry,'' \emph{Physics Reports} 492
(2010), 1--30. DOI: \href{https://doi.org/10.1016/j.physrep.2010.03.002}{10.1016/j.physrep.2010.03.002}.
\end{thebibliography}

\section*{Declarations}
\textbf{Scope.} \texttt{NO\_BAD\_EULER\_OR\_ROOT\_NUMBER}; no target
arithmetic or Hilbert--Polya claim.  \textbf{Data and code.} All rows and
audits are released with HCS-C224.  \textbf{AI-use disclosure.} Generative
tools assisted drafting and code generation; the internal artifact chain
checked displayed claims and metadata.  This is not external peer review.
\end{document}
